%!TEX root = ../../../main.tex

\chapter{Reverse Engineering: Habe ich verstanden, was gebaut wurde?}
\label{chap:reverse}

Den Code habe ich mit einem KI-Assistenten geschrieben. Damit ich trotzdem jedes Teil verstehe, bin ich nach jedem Meilenstein rückwärts vorgegangen. Ich habe nicht gefragt, was der Code tun soll, sondern nachgesehen, was er tut. Das geschah in drei Schritten. \emph{Erstens} steht in \texttt{docs/reverse\_engineering.md} für jede Datei, was hinein- und herausgeht, welche Bibliothek sie nutzt und welcher Test sie absichert. \emph{Zweitens} prüfen 104 Unit-Tests die Bausteine einzeln, etwa \enquote{eine neue Überschrift ersetzt die alte auf gleicher Ebene}. Dabei sperrt ein Test-Hook jeden Modellaufruf, sodass das Verhalten nachweislich aus dem Code kommt. \emph{Drittens} habe ich einzelne Fragen durch das ganze System verfolgt und jede Zwischendatei angesehen.

\paragraph{Eine Frage von Anfang bis Ende.} Frage T14 lautet: \enquote{Wurde die Schmutzwassergebühr in der Wasserabrechnung 2024 zu dem Satz berechnet, der im Wasservertrag steht?}
\begin{enumerate}[nosep, leftmargin=*]
    \item \emph{Parser:} PyMuPDF erkennt auf der Wasserrechnung nur die Kopfzeile der Tabelle. Der Parser baut die Zeilen aus den Wortpositionen nach (\texttt{"reconstructed": true}, Listing~\ref{lst:chunk}).
    \item \emph{Suche:} Die Zeile hat $\text{base}=0{,}77$. \enquote{Wasser} steht im Dateinamen, also $\text{boost}=1{,}5$ und $\text{final}=1{,}15$. Alle sechs Belege stammen aus der Rechnung, der Wasser\emph{vertrag} ist nicht dabei.
    \item \emph{Entwurf:} qwen2.5:7b liefert als JSON den Wert \enquote{2,80\,€/m³} aus Beleg~2, eine unnötige Rechnung \enquote{416,50} und die Antwort \enquote{Ja, \dots\ wie im Vertrag festgelegt [2]}.
    \item \emph{Prüfung und Zeit:} Code findet \enquote{2,80\,€/m³} in Beleg~2 und verwirft die unnötige Rechnung. Von 98\,s gehen 76\,s auf das Lesen des Prompts (1436 Tokens), 22\,s auf die Antwort.
\end{enumerate}
Daraus habe ich drei Dinge gelernt. Erstens kommt die Wartezeit vom langen Kontext, nicht von der Antwort. Zweitens stimmt die Zahl, aber \enquote{wie im Vertrag festgelegt} ist nicht belegt. Drittens sieht meine Zahlprüfung das nicht, weil sie nur Zahlen prüft. Deshalb nutze ich zusätzlich RAGAS-Faithfulness (Kapitel~\ref{chap:evaluation}).

\paragraph{Was dabei aufgefallen ist.} In der ersten Version (Version1) fand ich so Fehler, die kein Endergebnis verraten hätte. Der Parser stürzte bei allen PDF-Verträgen ab, weil eine Funktion in drei Fällen drei statt zwei Werte zurückgab. Der Kapitelpfad stapelte \enquote{§ 1 > § 2 > § 3}, statt § 1 durch § 2 zu ersetzen. Der Reranker lud beim Start ein Modell aus dem Internet, das System war also nicht offline. 15 Zusatzwerkzeuge, davon fünf für einen Wissensgraphen, waren nie gegen \enquote{ohne} gemessen worden; ich habe sie entfernt. Auch meine eigene Bewertung hatte einen Fehler: Sie las \enquote{wenn nicht gekündigt} als Verneinung und erkannte die neue Ablehnungsformel nicht. Version2 sah dadurch 11 bis 14 Prozentpunkte schlechter aus. Seitdem hat auch die Bewertung Tests.
